\documentclass{beamer}
\usepackage{amsmath}
\usepackage{amssymb}
\usepackage{graphicx}
\usepackage{algorithm,algorithmic}
\usepackage{listings}
\usepackage{color}

\definecolor{codegreen}{rgb}{0,0.6,0}
\definecolor{codegray}{rgb}{0.5,0.5,0.5}
\definecolor{codepurple}{rgb}{0.58,0,0.82}
\definecolor{backcolour}{rgb}{0.95,0.95,0.92}

\lstset{
   backgroundcolor=\color{backcolour},   
   commentstyle=\color{codegreen},
   keywordstyle=\color{magenta},
   numberstyle=\tiny\color{codegray},
   stringstyle=\color{codepurple},
   basicstyle=\footnotesize\ttfamily,
   breakatwhitespace=false,         
   breaklines=true,                 
   captionpos=b,                    
   keepspaces=true,                 
   numbers=left,                    
   numbersep=5pt,                  
   showspaces=false,                
   showstringspaces=false,
   showtabs=false,                  
   tabsize=2,
   language=Python
}

\title{Fast Fourier Transform (FFT) Algorithm}
\author{Morten HJ}
\date{Old notes from Comp Phys on FFT}

\usetheme{Madrid}

\begin{document}

\frame{\titlepage}

\section{Introduction}

\begin{frame}{Motivation}
  \begin{itemize}
  \item The Fourier Transform is a fundamental tool in signal processing, physics, engineering, and applied mathematics
  \item Direct computation of the Discrete Fourier Transform (DFT) is computationally expensive: $O(N^2)$
  \item The Fast Fourier Transform (FFT) reduces this to $O(N \log N)$
  \item Applications include:
    \begin{itemize}
    \item Signal processing (audio, image, video)
    \item Solving partial differential equations
    \item Polynomial multiplication
    \item Data compression
    \item Spectral analysis
    \end{itemize}
  \end{itemize}
\end{frame}

\begin{frame}{Historical Context}
  \begin{itemize}
  \item First discovered by Gauss in 1805 (unpublished) for calculating asteroid orbits
  \item Rediscovered by Cooley and Tukey in 1965
  \item Popularized with the advent of digital computers
  \item One of the most important numerical algorithms of the 20th century
  \end{itemize}
\end{frame}

\section{Mathematical Foundations}

\begin{frame}{Discrete Fourier Transform (DFT)}
  The DFT of a sequence $x[n]$ of length $N$ is defined as:
  \[ X[k] = \sum_{n=0}^{N-1} x[n] e^{-j 2\pi k n / N}, \quad k = 0,1,\ldots,N-1 \]
  where:
  \begin{itemize}
  \item $x[n]$ is the input sequence (time domain)
  \item $X[k]$ is the output sequence (frequency domain)
  \item $N$ is the number of samples
  \item $j = \sqrt{-1}$ is the imaginary unit
  \end{itemize}
\end{frame}

\begin{frame}{Inverse DFT}
  The inverse transform is given by:
  \[ x[n] = \frac{1}{N} \sum_{k=0}^{N-1} X[k] e^{j 2\pi k n / N}, \quad n = 0,1,\ldots,N-1 \]
  \begin{block}{Key Properties}
    \begin{itemize}
    \item Linearity: $a x_1[n] + b x_2[n] \leftrightarrow a X_1[k] + b X_2[k]$
    \item Periodicity: $X[k+N] = X[k]$
    \item Symmetry: For real $x[n]$, $X[N-k] = X^*[k]$
    \item Convolution: Multiplication in one domain equals convolution in the other
    \end{itemize}
  \end{block}
\end{frame}

\section{FFT Algorithm}

\begin{frame}{FFT Basic Idea}
  \begin{itemize}
  \item Exploits symmetry and periodicity of the complex exponential
  \item Decimates the DFT computation into smaller DFTs
  \item Most common approach: Cooley-Tukey algorithm (radix-2)
  \item Requires $N$ to be a power of 2 (can be generalized)
  \end{itemize}
  \begin{center}
%    \includegraphics[width=0.5\textwidth]{butterfly.png}
  \end{center}
\end{frame}

\begin{frame}{Divide and Conquer Approach}
  Split the DFT sum into even and odd terms:
  \begin{align*}
    X[k] &= \sum_{n=0}^{N-1} x[n] W_N^{kn} \\
    &= \sum_{m=0}^{N/2-1} x[2m] W_N^{k(2m)} + \sum_{m=0}^{N/2-1} x[2m+1] W_N^{k(2m+1)} \\
    &= \sum_{m=0}^{N/2-1} x_{\text{even}}[m] W_{N/2}^{km} + W_N^k \sum_{m=0}^{N/2-1} x_{\text{odd}}[m] W_{N/2}^{km} \\
    &= E[k] + W_N^k O[k]
  \end{align*}
  where $W_N = e^{-j 2\pi/N}$ is the twiddle factor.
\end{frame}

\begin{frame}{Recursive Structure}
  \begin{itemize}
  \item The DFT of size $N$ is decomposed into:
    \begin{itemize}
    \item Two DFTs of size $N/2$ (even and odd terms)
    \item $O(N)$ complex multiplications and additions
    \end{itemize}
  \item Applying this recursively gives $O(N \log N)$ complexity
  \end{itemize}
  \begin{center}
%    \includegraphics[width=0.7\textwidth]{fft_recursion.png}
  \end{center}
\end{frame}

\begin{frame}{Butterfly Operation}
  The basic computational unit of the FFT:
  \begin{center}
%    \includegraphics[width=0.5\textwidth]{butterfly_diagram.png}
  \end{center}
  Mathematically:
  \begin{align*}
    X[m] &= E[m] + W_N^m O[m] \\
    X[m+N/2] &= E[m] - W_N^m O[m]
  \end{align*}
\end{frame}

\section{Implementation}

\begin{frame}{Pseudocode for Radix-2 FFT}
  \begin{algorithm}[H]
    \begin{algorithmic}[1]
      \REQUIRE $N$ is a power of 2, $x$ is input array of size $N$
      \IF{$N == 1$}
      \RETURN $x$
      \ENDIF
      \STATE $x_{\text{even}} \gets [x[0], x[2], \ldots, x[N-2]]$
      \STATE $x_{\text{odd}} \gets [x[1], x[3], \ldots, x[N-1]]$
      \STATE $E \gets \text{FFT}(x_{\text{even}})$
      \STATE $O \gets \text{FFT}(x_{\text{odd}})$
      \STATE $W_N \gets e^{-2\pi j / N}$
      \STATE $W \gets 1$
      \FOR{$k = 0$ to $N/2 - 1$}
      \STATE $X[k] \gets E[k] + W \cdot O[k]$
      \STATE $X[k + N/2] \gets E[k] - W \cdot O[k]$
      \STATE $W \gets W \cdot W_N$
      \ENDFOR
      \RETURN $X$
    \end{algorithmic}
    \caption{Recursive Radix-2 FFT}
  \end{algorithm}
\end{frame}

\begin{frame}[fragile]{Python Implementation}
  \begin{lstlisting}
    import numpy as np

    def fft(x):
    N = len(x)
    if N == 1:
    return x
    even = fft(x[::2])
    odd = fft(x[1::2])
    T = [np.exp(-2j * np.pi * k / N) * odd[k]
      for k in range(N // 2)]
    return [even[k] + T[k] for k in range(N // 2)] + \
    [even[k] - T[k] for k in range(N // 2)]
  \end{lstlisting}
\end{frame}

\begin{frame}{Iterative Implementation}
  \begin{itemize}
  \item Recursive implementation has overhead
  \item More efficient: iterative version with bit-reversal permutation
  \item Three main steps:
    \begin{enumerate}
    \item Bit-reverse the input array
    \item Perform butterfly operations at each level
    \item Combine results
    \end{enumerate}
  \end{itemize}
\end{frame}

\section{Complexity Analysis}

\begin{frame}{Computational Complexity}
  \begin{itemize}
  \item Let $N = 2^m$
  \item At each level $k$ ($k = 1$ to $m$):
    \begin{itemize}
    \item $2^{k-1}$ butterfly operations
    \item Each butterfly has 1 complex multiplication and 2 additions
    \item Total operations per level: $N/2$ butterflies $\times$ 3 operations
    \end{itemize}
  \item Total levels: $\log_2 N$
  \item Total operations: $\frac{3}{2} N \log_2 N$
  \item Thus, complexity is $O(N \log N)$ vs. $O(N^2)$ for DFT
  \end{itemize}
\end{frame}

\begin{frame}{Comparison with Direct Computation}
  For $N = 1024$:
  \begin{itemize}
  \item Direct DFT: $N^2 = 1,048,576$ complex multiplications
  \item FFT: $\frac{N}{2} \log_2 N = 5120$ complex multiplications
  \item Speedup factor: ~205 times
  \end{itemize}
  \begin{center}
    \begin{tabular}{|c|c|c|}
      \hline
      $N$ & DFT Operations & FFT Operations \\
      \hline
      32 & 1,024 & 80 \\
      64 & 4,096 & 192 \\
      128 & 16,384 & 448 \\
      256 & 65,536 & 1,024 \\
      512 & 262,144 & 2,304 \\
      1024 & 1,048,576 & 5,120 \\
      \hline
    \end{tabular}
  \end{center}
\end{frame}

\section{Variations and Extensions}

\begin{frame}{Different FFT Algorithms}
  \begin{itemize}
  \item \textbf{Radix-2}: Most common, requires $N$ power of 2
  \item \textbf{Radix-4}: More efficient for certain architectures
  \item \textbf{Mixed-radix}: For arbitrary $N$ (prime factor algorithm)
  \item \textbf{Bluestein's algorithm}: For arbitrary $N$ (Chirp-Z transform)
  \item \textbf{Winograd FFT}: Minimizes multiplications
  \end{itemize}
\end{frame}

\begin{frame}{Multidimensional FFT}
  \begin{itemize}
  \item Separable transform: Apply 1D FFT along each dimension
  \item For 2D image of size $M \times N$:
    \[ X[k,l] = \sum_{m=0}^{M-1} \sum_{n=0}^{N-1} x[m,n] e^{-j2\pi (km/M + ln/N)} \]
  \item Computed as:
    \begin{enumerate}
    \item Apply 1D FFT to each row ($M \times O(N \log N)$)
    \item Apply 1D FFT to each column ($N \times O(M \log M)$)
    \end{enumerate}
  \item Total complexity: $O(MN \log MN)$
  \end{itemize}
\end{frame}

\section{Applications}

\begin{frame}{Practical Applications}
  \begin{itemize}
  \item \textbf{Signal Processing}:
    \begin{itemize}
    \item Filtering (convolution via multiplication in frequency domain)
    \item Spectral analysis
    \item Audio compression (MP3, AAC)
    \end{itemize}
  \item \textbf{Image Processing}:
    \begin{itemize}
    \item JPEG compression (2D FFT on 8x8 blocks)
    \item Filtering and enhancement
    \item Feature extraction
    \end{itemize}
  \item \textbf{Numerical Analysis}:
    \begin{itemize}
    \item Solving PDEs (spectral methods)
    \item Fast polynomial multiplication
    \end{itemize}
  \item \textbf{Other Fields}:
    \begin{itemize}
    \item Quantum computing (QFT)
%      \option{Medical imaging (MRI reconstruction)
      \item Astronomy (interferometry)
    \end{itemize}
  \end{itemize}
\end{frame}

\begin{frame}{Example: Polynomial Multiplication}
  \begin{itemize}
  \item Multiply two polynomials $A(x)$ and $B(x)$ of degree $n-1$
  \item Naive approach: $O(n^2)$ operations
  \item Using FFT:
    \begin{enumerate}
    \item Evaluate $A$ and $B$ at $2n$ points using FFT ($O(n \log n)$)
    \item Pointwise multiply the results ($O(n)$)
    \item Interpolate to get coefficients using inverse FFT ($O(n \log n)$)
    \end{enumerate}
  \item Total complexity: $O(n \log n)$
  \end{itemize}
\end{frame}

\section{Conclusion}

\begin{frame}{Summary}
  \begin{itemize}
  \item FFT is an efficient algorithm to compute the DFT
  \item Reduces complexity from $O(N^2)$ to $O(N \log N)$
  \item Based on divide-and-conquer strategy
  \item Numerous applications across science and engineering
  \item Modern implementations (FFTW) are highly optimized
  \end{itemize}
\end{frame}

\begin{frame}{Further Reading}
  \begin{itemize}
  \item Cooley, J. W.; Tukey, J. W. (1965). ``An algorithm for the machine calculation of complex Fourier series''. Math. Comput. 19: 297-301.
  \item Brigham, E. O. (1988). \textit{The Fast Fourier Transform and Its Applications}. Prentice-Hall.
  \item Van Loan, C. (1992). \textit{Computational Frameworks for the Fast Fourier Transform}. SIAM.
  \item FFTW: http://www.fftw.org/
  \item Numerical Recipes in C: The Art of Scientific Computing
  \end{itemize}
\end{frame}

\end{document}
